Cancerul de piele, în special melanomul, este foarte tratabil atunci când este
depistat precoce, însă prognosticul se înrăutățește dramatic în stadiile
avansate; totodată, accesul la îngrijire dermatologică de specialitate este
distribuit inegal la nivel mondial. Această lucrare prezintă DermAI, un sistem
end-to-end de analiză a leziunilor cutanate bazat pe inteligență artificială,
care sprijină screeningul precoce printr-un instrument disponibil atât pe
telefoane inteligente, cât și în browsere web.

Partea de cercetare a lucrării antrenează un model Vision Transformer (ViT) pe
setul de date ISIC 2024 pentru clasificarea leziunilor cutanate în benigne și
maligne, acordând o atenție deosebită gestionării dezechilibrului sever al
claselor. Modelul este comparat cu o rețea convoluțională ResNet-18, antrenată în
condiții identice, și cu un model lingvistic de mari dimensiuni Gemini în regim
zero-shot, folosind metrica oficială ISIC 2024, completată cu teste de
semnificație statistică și analiză de calibrare. Partea inginerească extinde
modelul antrenat într-o aplicație completă, cuprinzând clienți mobili și web, un
strat MCP care servește modelul, validarea imaginilor și explicații în limbaj
natural bazate pe Gemini, precum și stocare în cloud autentificată și separată pe
utilizatori.

\textbf{Cuvinte cheie}: cancer de piele, melanom, învățare profundă, Vision
Transformer, ISIC 2024, dezechilibru al claselor, model lingvistic de mari
dimensiuni, sănătate mobilă
